% !TEX program = xelatex

\documentclass[11pt,a4paper]{ctexart}

% ====== 字体 ======
\usepackage{fontspec}
\setmainfont{Latin Modern Roman}
\usepackage{unicode-math}
% 修改：XITS Math 拥有原生粗体数学字形，解决 \symbfit 加粗无效的问题
\setmathfont{XITS Math}

% ====== 颜色支持 ======
\usepackage{xcolor}

% ====== 页面设置 ======
\usepackage[margin=2.5cm]{geometry}
\usepackage{setspace}
\onehalfspacing
\pagestyle{plain}

% ====== 数学宏包 ======
\usepackage{mathtools}
\usepackage{amsmath}
% \usepackage{amssymb}   % 删除
% \usepackage{bm}        % 彻底删除

% ====== 三线表 + 数字对齐 ======
\usepackage{booktabs}
\usepackage{siunitx}
\usepackage{enumitem}

% ====== 交叉引用与超链接 ======
\usepackage[colorlinks=true,
            linkcolor=blue,
            citecolor=green,
            urlcolor=cyan]{hyperref}
\usepackage[noabbrev,nameinlink]{cleveref}

% ====== 框架图形 ======
\usepackage{tcolorbox}  % 引入 tcolorbox 宏包
\tcbuselibrary{skins, breakable} % 加载所需库，以实现更多样式和自动分页

% 定义一个名为 block 的环境
\newtcolorbox{block}{
    colback=white,        % 背景色
    colframe=black!75,    % 边框颜色
    boxrule=0.5mm,        % 边框粗细
    arc=2mm,              % 圆角大小
    boxsep=5pt,           % 边框与内容的间距
    left=6pt, right=6pt,  % 左右内边距
    top=6pt, bottom=6pt,  % 上下内边距
    breakable             % 允许盒子在跨页时自动断开
}

% ====== 自定义命令 ======
\newcommand{\dd}{\mathrm{d}}
\newcommand{\source}[1]{\par \hfill ——#1}

\title{0阶 Bessel 函数的 Laplace 变换}
\author{王晨暄}
\date{\today}

\begin{document}
\maketitle

\begin{block}

我们知道对于 Bessel 函数 $J_0(t)$，其 Laplace 变换为：

\begin{equation}
    \mathcal{L}\{J_0(t)\} = \int_0^\infty e^{-st} J_0(t) \dd{t} = \frac{1}{\sqrt{s^2 + 1}}
    \label{eq:bessel_laplace}
\end{equation}

公式\eqref{eq:bessel_laplace}说明 $J_0(t)$ 的 Laplace 变换
的结果为 $\dfrac{1}{\sqrt{s^2 + 1}}$，$\dfrac{1}{\sqrt{s^2 + 1}}$的 Laplace 反演为 $J_0(t)$。
这个结果在信号处理和控制理论中有广泛的应用。
下面我们从物理问题出发，使用特定电荷分布与边界条件下的静电场来推导这个公式。

\end{block}

\bigskip

{\kaishu
    假设无限大导体平面$z=0$接地，空间中的电荷分布为$\rho(\symbfit{r})=q\delta(z-z_0)\delta(x)\delta(y)$。
    利用电象法求解$z\geq 0$的区域内的电势分布。同时考虑 Poisson 方程的积分变换，导出方程\eqref{eq:bessel_laplace}。
}

\bigskip

首先根据电象法可知，系统的空间电势函数为：

\begin{equation}
    \phi(\symbfit{r}) = \frac{q}{4\pi\epsilon_0}\left(\frac{1}{|\symbfit{r}-\symbfit{r}_0|} - \frac{1}{|\symbfit{r}-\symbfit{r}_0'|}\right)
    \label{eq:potential}
\end{equation}

即：

\begin{equation}
    \phi(\symbfit{r}) = \frac{q}{4\pi\epsilon_0}\left(\frac{1}{\sqrt{(z-z_0)^2 + r^2}} - \frac{1}{\sqrt{(z+z_0)^2 + r^2}}\right)
    \label{eq:potential_explicit}
\end{equation}

由于在上半空间中电势函数满足 Poisson 方程：

\begin{equation}
    \nabla^2 \phi(\symbfit{r}) = -\frac{\rho(\symbfit{r})}{\epsilon_0}
\end{equation}

得到解：

\begin{equation}
    \phi(\symbfit{r}) = \frac{q}{4\pi\epsilon_0}\left(\frac{1}{\sqrt{(z-z_0)^2 + r^2}}\right) + \int_{0}^{\infty}A(k)J_0(kr)e^{-kz}\dd{k}
    \label{eq:potential_solution}
\end{equation}

其中第二项为满足边界条件的齐次解。对比\eqref{eq:potential_explicit}与\eqref{eq:potential_solution}，我们可以得到：

\begin{equation}
    \int_{0}^{\infty}A(k)J_0(kr)e^{-kz}\dd{k} = -\frac{q}{4\pi\epsilon_0}\left(\frac{1}{\sqrt{(z+z_0)^2 + r^2}}\right)
    \label{eq:integral_equation}
\end{equation}

注意到$J_0(0)=1$，因此在上述方程中取$r=0$，可以解出$A(k)$:

\begin{equation}
    A(k) = \mathcal{L}^{-1}\left(-\frac{q}{4\pi\epsilon_0} \frac{1}{\left|z+z_0\right|}\right)
\end{equation}

根据 Laplace 变换的性质，我们可以得到：

\begin{equation}
    A(k)= -\frac{q}{4\pi\epsilon_0} e^{-kz_0}
\end{equation}

代入\eqref{eq:integral_equation}，化简后得到：

\begin{equation}
    \mathcal{L}\{J_0(t)\} = \int_0^\infty e^{-st} J_0(t) \dd{t} = \frac{1}{\sqrt{s^2 + 1}}
\end{equation}

本文旨在提供一个启发式的物理类比（heuristic physical analogy），用于辅助记忆该积分变换公式，而非严格的数学证明。

\end{document}